\documentclass[it]{../../elegantbook}
% ========== Copertina e Informazioni Base ==========
\title{Documento di test ElegantBook}
\subtitle{Caratteristiche principali del modello}
\author{Ingegnere di test}
\institute{Gruppo di test del modello}
\date{\today}
\version{v1.0}
\extrainfo{Questo documento è solo per il test delle funzioni del modello, non ha alcun valore accademico}
\bioinfo{Lotto di test}{TEST-FULL-20260506}
\logo{../../figure/logo-blue.png}
\cover{../../figure/cover.jpg}

% ========== Indice e Impostazioni Base ==========
\setcounter{tocdepth}{3} % Indice visibile fino al livello 3

% ========== Bibliografia di test ==========
\addbibresource[location=local]{../en/reference.bib}

\begin{document}

\maketitle
\frontmatter
\tableofcontents

% ========== Prefazione (con sezioni) ==========
\pagenumbering{Roman}
\setcounter{page}{0}
\chapter*{Prefazione di test}
\markboth{Prefazione di test}{}
\section*{Descrizione documento}
\markright{Descrizione documento}
Questo documento è un caso di test completo per ElegantBook, copre tutte le funzionalità base e avanzate del modello. Tutti i contenuti sono dati di test fittizi senza significato pratico.
\newpage
\section*{Ambito di test}
\markright{Ambito di test}
Questo test include: copertina, indice, caratteri, elenchi multilivello, formule complesse, tutti gli ambienti matematici, riferimenti incrociati, figure e tabelle, note marginali, riassunti di capitolo, esercizi, bibliografia, appendici.

\mainmatter

% ========== Capitolo 1: Impaginazione base ed elenchi multilivello ==========
\chapter{Test impaginazione base ed elenchi}
\begin{introduction}[Elementi di test di questo capitolo]
\item Caratteri cinesi e formattazione testo
\item Annidamento elenco non ordinato a 3 livelli
\item Annidamento elenco ordinato a 3 livelli
\item Paragrafi normali e testo enfatizzato
\end{introduction}

\section{Test caratteri e testo}
Testo normale del corpo. \textbf{Grassetto}, \textit{Corsivo}, \underline{Sottolineato}, \textcolor{red}{Testo rosso}.

\section{Test elenco non ordinato a 3 livelli}
\begin{itemize}
\item Elemento di test livello 1 A
\item Elemento di test livello 1 B
    \begin{itemize}
    \item Elemento di test livello 2 B1
    \item Elemento di test livello 2 B2
        \begin{itemize}
        \item Elemento di test livello 3 B2-1
        \item Elemento di test livello 3 B2-2
            \begin{itemize}
            \item Elemento di test livello 4 B2-1a
            \item Elemento di test livello 4 B2-2b
            \end{itemize}
        \end{itemize}
    \end{itemize}
\end{itemize}

\section{Test elenco ordinato a 3 livelli}
\begin{enumerate}
\item Test ordinato livello 1 1
\item Test ordinato livello 1 2
    \begin{enumerate}
    \item Test ordinato livello 2 2-1
    \item Test ordinato livello 2 2-2
        \begin{enumerate}
        \item Test ordinato livello 3 2-2-1
        \item Test ordinato livello 3 2-2-2
            \begin{enumerate}
            \item Test ordinato livello 4 2-2-2-1
            \item Test ordinato livello 4 2-2-2-2
            \end{enumerate}
        \end{enumerate}
    \end{enumerate}
\end{enumerate}

% ========== Capitolo 2: Test formule matematiche complesse ==========
\chapter{Test formule matematiche complesse}
\begin{introduction}
\item Formula complessa a riga singola
\item Formule multilinea (align)
\item Matrici / Determinanti
\item Integrali multipli / Somme / Limiti / Frazioni
\item Etichette formule e riferimenti incrociati
\end{introduction}

\section{Formula complessa a riga singola}
Formule complesse in linea: $\sum_{i=1}^n \frac{x_i^2}{\sqrt{y_i+1}} \geq \bar{x}^2$, $\lim_{n\to\infty} \left(1+\frac{1}{n}\right)^n = e$.

\section{Formule multilinea e matrici}
\begin{align}
a^2 + b^2 &= c^2 \label{eq:gougu} \\
\int_{0}^{1}\int_{0}^{1} f(x,y) dxdy &= 1 \label{eq:double} \\
\begin{vmatrix}
a & b \\
c & d
\end{vmatrix} &= ad-bc \label{eq:det}
\end{align}

\section{Formule con operatori grandi}
\begin{equation}\label{eq:series}
\sum_{k=0}^{\infty} \frac{x^k}{k!} = e^x,\quad \oint_{C} Pdx+Qdy = \iint_{D}\left(\frac{\partial Q}{\partial x}-\frac{\partial P}{\partial y}\right)dxdy
\end{equation}

Riferimento formula: Teorema di Pitagora \eqref{eq:gougu}, integrale doppio \eqref{eq:double}, determinante \eqref{eq:det}, serie \eqref{eq:series}.

% ========== Capitolo 3: Test tutti gli ambienti matematici ==========
\chapter{Test completo ambienti matematici}
\begin{introduction}
\item Definizione / Teorema / Lemma / Corollario / Assioma / Postulato
\item Proposizione / Esempio / Problema / Esercizio
\item Nota / Osservazione / Conclusione / Ipotesi / Proprietà
\item Ambiente soluzione / dimostrazione
\end{introduction}

\section{Test ambienti}

% 1. Definizione
\ifdefined\simpleTest
\begin{definition}[Definizione di test]\label{def:test}
\else
\begin{definition}[Definizione di test]{test}
\fi
Questo è il contenuto di test dell'ambiente definizione, senza significato matematico reale, solo per verifica stile.
\end{definition}

% 2. Teorema
\ifdefined\simpleTest
\begin{theorem}[Teorema di test]\label{thm:test}
\else
\begin{theorem}[Teorema di test]{test}
\fi
Se $x>0$, allora $x+1>1$, corrispondente alla formula \eqref{eq:gougu}.
\end{theorem}

% 3. Lemma
\ifdefined\simpleTest
\begin{lemma}[Lemma di test]\label{lem:test}
\else
\begin{lemma}[Lemma di test]{test}
\fi
Per ogni numero reale $x$, vale $x^2\geq0$.
\end{lemma}

% 4. Corollario
\ifdefined\simpleTest
\begin{corollary}[Corollario di test]\label{cor:test}
\else
\begin{corollary}[Corollario di test]{test}
\fi
Dal Lemma \ref{lem:test} si ottiene: $(x-1)^2\geq0$.
\end{corollary}

% 5. Assioma
\ifdefined\simpleTest
\begin{axiom}[Assioma di test]\label{axi:test}
\else
\begin{axiom}[Assioma di test]{test}
\fi
Assioma di test, senza significato pratico.
\end{axiom}

% 6. Postulato
\ifdefined\simpleTest
\begin{postulate}[Postulato di test]\label{pos:test}
\else
\begin{postulate}[Postulato di test]{test}
\fi
Postulato di test, solo per verifica ambiente.
\end{postulate}

% 7. Proposizione
\ifdefined\simpleTest
\begin{proposition}[Proposizione di test]\label{pro:test}
\else
\begin{proposition}[Proposizione di test]{test}
\fi
Contenuto proposizione di test, collegato alla Definizione \ref{def:test}.
\end{proposition}

% 8. Esempio
\begin{example}\label{exa:test}
Test ambiente esempio, numerazione automatica.
\end{example}

% 9. Problema
\begin{problem}\label{probl:test}
Test ambiente problema, per proporre quesiti da risolvere.
\end{problem}

% 10. Esercizio
\begin{exercise}\label{exer:test}
Calcola: $1+2+3+\dots+10$.
\end{exercise}

% 11. Nota
\begin{note}
Ambiente nota: senza numerazione, per avvertenze chiave.
\end{note}

% 12. Osservazione
\begin{remark}
Ambiente osservazione: per spiegazioni supplementari.
\end{remark}

% 13. Conclusione
\begin{conclusion}
Ambiente conclusione: contenuto riassuntivo di test.
\end{conclusion}

% 14. Ipotesi
\begin{assumption}
Ambiente ipotesi: impostazione condizioni di test.
\end{assumption}

% 15. Proprietà
\begin{property}
Ambiente proprietà: caratteristiche oggetto di test.
\end{property}

% 16. Soluzione
\begin{solution}
Soluzione Esercizio \ref{exer:test}: La somma vale $55$ (risposta di test).
\end{solution}

% 17. Dimostrazione
\begin{proof}
Dimostrazione Teorema \ref{thm:test}: Poiché $x>0$, sommando 1 a entrambi i lati si dimostra la tesi.
\end{proof}

% ========== Capitolo 4: Riferimenti incrociati · Figure · Note marginali ==========
\chapter{Riferimenti incrociati e funzionalità avanzate}
\begin{introduction}
\item Tutti i tipi di riferimenti incrociati
\item Test semplici figure e tabelle
\item Test funzionalità note marginali
\item Verifica navigazione capitoli
\end{introduction}

\section{Riepilogo riferimenti incrociati}
Definizione \ref{def:test}, Teorema \ref{thm:test}, Lemma \ref{lem:test}, Corollario \ref{cor:test}, Assioma \ref{axi:test}, Postulato \ref{pos:test}, Proposizione \ref{pro:test}, Esempio \ref{exa:test}, Problema \ref{probl:test}, Esercizio \ref{exer:test}.

\section{Test figure e tabelle}
\begin{figure}[h]
    \centering
    \includegraphics[width=0.5\textwidth]{../../image/scatter.jpg}
    \caption{Figura di test}\label{fig:test}
\end{figure}
\begin{table}[h]
    \centering
    \begin{tabular}{c*{6}{c}}
    \toprule
        $n$ & 1 & 2 & 3 & 4 & 5 & 6 \\
    \midrule
        $a_n$ & 1 & 1 & 2 & 3 & 5 & 8 \\ 
    \bottomrule
    \end{tabular}
    \caption{Tabella di test}\label{tab:test}
\end{table}
Riferimento figure/tabelle: \figref{fig:test} è la figura di test, \tabref{tab:test} è la tabella di test.

% ========== Esercizi fine capitolo ==========
\begin{problemset}[Esercizi di test del capitolo]
\item Verifica lo stile visualizzazione della Definizione \ref{def:test}
\item Deriva la formula \eqref{eq:double}
\item Controlla se la Figura \ref{fig:test} viene visualizzata correttamente
\end{problemset}

% ========== Bibliografia ==========
\nocite{*}
\printbibliography[heading=bibintoc]

% ========== Appendice ==========
\appendix
\chapter{Appendice di test}
Il contenuto dell'appendice serve a verificare capitoli appendice, intestazioni, numerazione pagine e stile impaginazione, senza informazioni reali.
\section{Test Appendice 1}
Il contenuto dell'appendice serve a verificare capitoli appendice, intestazioni, numerazione pagine e stile impaginazione, senza informazioni reali.
\newpage
\section{Test Appendice 2}
Il contenuto dell'appendice serve a verificare capitoli appendice, intestazioni, numerazione pagine e stile impaginazione, senza informazioni reali.
\end{document}